\documentclass{standalone}
\usepackage[european]{circuitikz}

\begin{document}
	\begin{circuitikz} 
		\draw
		% Transistor Q1
		(0,0) node[npn, anchor=E] (Q1) {}
		(Q1.B) to[R, l_=${R_B=1k\Omega}$] ++(-2, 0)
		to[V, l_=$U_1$] ++(0,-2)
		node[rground]{} coordinate (gndL)
		(Q1.E) -- (Q1.E |- gndL)
		node[rground]{} coordinate (gndC)
		(Q1.C) to[R, l=${R_C=22\Omega}$] ++(0, 2) coordinate(Rctop)
		to[short] ++(5, 0) coordinate (a)
		to[sV, v=$0V-10V$, l_=$U_2$] (a |- gndL)
		node[rground]{} coordinate (gndC)
		;
		
		\draw[draw=red, line width=1pt] 
		(Q1.C) to[short, *-] ++(2, 0) coordinate (K1a)
		to[voltmeter, color=red] (K1a|-gndL) coordinate (K1b)
		node[above, color=red, xshift=.4cm, rotate=-90] at ($ (K1a)!0.5!(K1b) $) {Kanal 1}
		(K1b) node[rground]{}
		;
		
		\draw[draw=blue, line width=1pt] 
		(Q1.C) ++(0, .2) coordinate(x)
		to[short, *-] ++(2, 0) coordinate (K2a)
		to[voltmeter, color=blue, -*] (K2a|-Rctop) coordinate (K2b)
		node[above, color=blue, xshift=.4cm, rotate=-90] at ($ (K2a)!0.5!(K2b) $) {Kanal 2}
		;
		
	\end{circuitikz}
	
	
\end{document}
